\section{Introduction}

LLM agents increasingly act through tools: they read private files, modify software, manage communication, and interact with external services. Their safety therefore depends not only on generated responses but also on the consequences of execution. Runtime safety frameworks such as AgentSpec~\cite{wang2026agentspec} allow developers to enforce interpretable safety rules independently of an agent's decisions. However, these frameworks assume that domain-experts can specify the conditions under which actions are safe. In practice, safety requirements are often expressed as high-level natural-language policies, while their enforcement requires precise, executable rules over agent actions and execution context. Translating such policies into enforceable contracts remains a challenging and largely manual process. 

Manually specified contracts may overlook subtle behavioral conditions, creating gaps that agents can inadvertently exploit or deliberately circumvent while remaining within the monitor's explicit rules~\cite{kale2026monitoring}. During internal evaluations in 2026, OpenAI models escaped network isolation to compromise Hugging Face infrastructure and harvest production credentials~\cite{openai2026huggingface}, and accessed several Australian government systems without authorization, retrieving internal files and credentials~\cite{openai2026australia}.

Consider a document assistant governed by a policy stating that confidential reports must only be accessible to authorized recipients. A seemingly reasonable rule checks recipient authorization before sharing a file. Yet this rule permits the assistant to upload a confidential report to a publicly accessible location and send its URL to an authorized board member. Although the intended recipient is authorized, the report is now accessible to anyone possessing the link. Conversely, prohibiting all public links would unnecessarily prevent the assistant from distributing public posters. The correct safety contract must distinguish confidential from public content, authorized from unauthorized audiences, and restricted from publicly accessible sharing mechanisms. Crucially, these distinctions may depend on the history of preceding operations rather than the final action alone. A monitor can enforce every specified rule correctly and still permit harm if the contract omits such conditions.
% The challenge is therefore not simply enforcing safety rules, but automatically discovering the behavioral distinctions that make those rules correct and generalizable. 

Existing approaches address related aspects of policy-driven safety. PolicyBank updates agents' policy understanding through structured memory and corrective feedback, but the acquired knowledge is not an independently enforced runtime contract~\cite{choi2026policybank}. AutoSpec refines programmable safety rules from labeled examples using inductive synthesis and predicate selection~\cite{ma2026autospec}. SEVerA evolves agent programs under supplied hard constraints, rather than synthesizing an independent safety contract~\cite{banerjee2026severa}. AgentSpec provides an interface for deploying customizable trigger--check--enforce rules, but relies on developers to supply suitable rule conditions~\cite{wang2026agentspec}. These approaches motivate a complementary problem: \emph{given a target agent, a high-level safety policy, and a dataset of agent trajectories labeled safe or unsafe, how can we automatically learn an interpretable behavioral model that can be enforced as a safety contract, without requiring manually specified rules?}

\paragraph{What $L^*$ can do.}
A safety contract of this kind is a \emph{prefix-closed language} over policy-relevant events: an execution prefix is accepted as long as no violation has occurred. Its minimal deterministic finite automaton (DFA) has one state per class of \emph{execution histories with the same permissible continuations}, such as ``a confidential report is staged'' versus ``the staged file is public''. These states are exactly what manual rules miss. Angluin's $L^*$~\cite{angluin1987learning} learns such a DFA without enumerating executions. It asks \emph{membership queries}, whether a given event sequence is acceptable, and \emph{equivalence queries}, whether a candidate DFA is correct and, if not, for a counterexample. Each counterexample reveals a missing history distinction and adds a state. With exact answers to both kinds of query, $L^*$ returns exactly the minimal DFA after at most $N-1$ counterexamples, where $N$ is its number of states, using polynomially many queries.

\paragraph{What stands in the way.}
$L^*$ assumes three things that an agent setting does not provide. What is safe is ultimately decided by the world, human judgment or the actual effects of actions, which is observable only on recorded trajectories, while written policies are incomplete summaries of that experience.
\textbf{(C1) A fixed alphabet.} $L^*$ learns over given symbols, but an agent emits concrete tool calls whose policy relevance depends on arguments, content and context. The symbols must be defined by predicates over agent observations, and they must preserve the ground-truth judgments: two executions mapped to the same word must be judged alike.
\textbf{(C2) A membership oracle.} $L^*$ needs a teacher that labels any event sequence, but ground truth exists only for the sequences of recorded trajectories, and $L^*$ asks about many more. An LLM that reads a policy can answer the rest, but $L^*$ reproduces every wrong answer exactly: in our study, one wrong label among 399 grew a four-state automaton to ten states.
\textbf{(C3) An equivalence oracle.} Nothing can confirm that a DFA is correct, and this query is what makes $L^*$'s result correct. It cannot be delegated to an LLM: an LLM that could certify a candidate automaton could simply write the automaton, and the result would rest on the judgment it is meant to check.

\paragraph{Our approach.}
We propose \system{}, which learns an enforceable AgentSpec contract from an instrumentable agent $A$, a dataset $\mathcal{D}$ of agent trajectories with ground-truth safety labels, such as an existing agent-safety benchmark, and a natural-language policy $P$ that summarizes this experience. The contract should reproduce the ground truth; the policy is an aid, not the specification. For (C1), an expert defines the event alphabet through predicates of any implementation, symbolic or LLM-based. For (C2), \system splits the teacher: the labels answer every membership query they determine, and an LLM that reads $P$ answers the rest, so the policy generalizes the experience to behavior the data does not show. For (C3), the dataset answers equivalence queries: \system replays $\mathcal{D}$ against the current DFA, and every disagreement is a counterexample. The LLM thus answers only local questions the data cannot settle, while the global question, whether the DFA is right, is answered by real feedback from the world. Where the labels determine a word, \system also records the LLM's answer, and disagreements show where the policy fails to summarize the experience. Finally, \system compiles the DFA into AgentSpec: event mappings provide triggers, DFA transitions implement the stateful check, and pre-effect mediation enforces it.

\paragraph{Guarantee.}
Whatever the LLM answers, the number of counterexamples is bounded by the number of distinct sequences in the dataset. If the labels are correct and the LLM is correct where they are silent, \system returns a DFA that agrees with every labeled trajectory after at most $N-1$ counterexamples, and it is never larger than the target's minimal DFA. If the dataset is also complete, the DFA agrees with the ground truth on all realizable behavior; if, in addition, the predicates preserve the ground-truth judgments and interception is faithful, the compiled contract blocks exactly the unacceptable actions (Section~\ref{sec:problem}). Existing datasets are not complete, so we measure generalization on held-out labeled trajectories, and we check each assumption empirically.

In summary, this work makes three contributions: (i) a formalization of agent safety-contract learning as active inference of a prefix-closed regular language over policy-relevant event predicates of any implementation, with explicit conditions under which $L^*$'s guarantees carry over when ground truth is available only on recorded trajectories; (ii) a split membership teacher, in which ground-truth labels answer where they apply and a policy-reading LLM generalizes beyond them; and (iii) a dataset-based equivalence check, which tests hypotheses against ground-truth-labeled agent trajectories instead of LLM judgments, with a bound on the number of counterexamples that holds even when the LLM errs. We compile the learned models into enforceable AgentSpec contracts and evaluate their safety, preservation of legitimate functionality, oracle cost, and runtime enforcement fidelity.
